% Ręcznie tworzona lista symboli - substytut w przypadku kompilowania za pomocą Overleaf
\flushleft

\begin{itemize}[noitemsep,topsep=0pt,parsep=0pt,partopsep=0pt,labelwidth=3cm,align=left,leftmargin=4cm]
    \item[API] interfejs programistyczny aplikacji (ang. Application Programming Interface)
    \item[BBOX] prostokąt otaczający obiekt (ang. Bounding Box)
    \item[CNN] splotowa sieć neuronowa (ang. Convolutional Neural Network)
    \item[CUDA] architektura procesorów firmy NVIDIA do obliczeń równoległych (ang. Compute Unified Device Architecture)
    \item[DDP] rozproszone uczenie równoległe danych (ang. Distributed Data Parallel)
    \item[DINO] detektor obiektów oparty na transformatorach (ang. DETR with Improved Denoising Anchor Boxes)
    \item[FPS] liczba klatek na sekundę (ang. Frames Per Second)
    \item[GPU] procesor graficzny (ang. Graphical Processing Unit)
    \item[IoU] miara przecięcia nad sumą (ang. Intersection over Union)
    \item[mAP] średnia precyzja dla wielu klas (ang. Mean Average Precision)
    \item[ML] uczenie maszynowe (ang. Machine Learning)
    \item[NMS] algorytm tłumienia niemaksymalnego (ang. Non-Maximum Suppression)
    \item[ONNX] otwarty format wymiany modeli neuronowych (ang. Open Neural Network Exchange)
    \item[SAM] model segmentacji dowolnych obiektów (ang. Segment Anything Model)
    \item[TFLite] zoptymalizowany format mobilny TensorFlow (ang. TensorFlow Lite)
    \item[VRAM] pamięć podręczna karty graficznej (ang. Video Random Access Memory)
    \item[YOLO] architektura detekcji obiektów w jednym przebiegu (ang. You Only Look Once)
\end{itemize}
